% macros for formatting references within the document
\newcommand{\fref}[1]     {Fig.~\ref{#1}}
\newcommand{\tref}[1]     {Tab.~\ref{#1}}
\newcommand{\eref}[1]     {Equation~(\ref{#1})} % equation
\newcommand{\aref}[1]     {Alg.~\ref{#1}} % algorithm
\newcommand{\apref}[1]     {Appendix~\ref{#1}} % appendix
\newcommand{\sref}[1]     {Section~\ref{#1}} % section
\newcommand{\pref}[1]   {Part~\ref{#1}}
\newcommand{\trefs}[2] {Tables~\ref{#1} and~\ref{#2}}

% Text macros
\newcommand{\eg}    {\textit{e.g.}}
\newcommand{\ie}    {\textit{i.e.}}
\newcommand{\todo}[1]  {\textcolor{red}{TODO: #1}}
\newcommand{\quotes}[1]{`#1'}
\newcommand{\dquotes}[1]{``#1''}

%% Math formatting - these are used by later macros 
\newcommand{\xmath}[1]	{\ensuremath{#1}\xspace}% math with a space if needed
\newcommand{\bmath}[1]	{\xmath{\bm{#1}}}	% this is the best math bold!

%% Basic math commands
% fractions 
\newcommand{\onehalf}	{\xmath{\dfrac{1}{2}}}
\newcommand{\onehalft}{\xmath{\tfrac{1}{2}}} % 1/2 formatted for in-text
\newcommand{\onequarter}	{\xmath{\dfrac{1}{4}}}
\newcommand{\onethird}	{\xmath{\dfrac{1}{3}}}
\newcommand{\twothirds}	{\xmath{\dfrac{2}{3}}}
\newcommand{\paren}[1]{\xmath{\left(#1\right)}}

% operators 
\newcommand{\conv} 		{\circledast}  % convolution symbol
\newcommand{\abs}[1]   {\xmath{\lvert #1 \rvert}} % | . | absolute value 
\newcommand{\norm}[1]   {\xmath{\left\lVert#1\right\rVert}} % || . ||
\newcommand{\normsq}[1]   {\xmath{\left\lVert#1\right\rVert^2}} % || . ||^2
\newcommand{\normrsq}[1]{\xmath{\| #1 \|^2}} % "regular" - not big
\newcommand{\real}[1]{\;\xmath{\text{real} \left\{ #1 \right\}) }}
\newcommand{\sign}[1]{\xmath{\mathrm{sign}\left(#1\right)}}
\renewcommand{\log}[1] {\xmath{\operatorname{log}\mleft( #1 \mright)}} % trick to get proper spacing all around

% define other helpful math stuff 
\newcommand{\by}      {\xmath{\times}} % the x signal in a dimension, e.g., m x n 
\renewcommand{\neg}         {\xmath{\text{-}}} % this negative sign sometimes looks better than usting '-' in math mode

% Common part names
\newcommand{\R}[1]      {\xmath{R_{#1}}} % for example, \R{1} produces R_1 and \R{2} produces R_2
\newcommand{\C}[1]      {\xmath{C_{#1}}} 

% Common signal names 
\newcommand{\Vout}  {\xmath{V_\text{out}}}
\newcommand{\Vin}  {\xmath{V_\text{in}}}
\newcommand{\Vleft}  {\xmath{V_\text{left}}}
\newcommand{\Vright}  {\xmath{V_\text{right}}}



% helpful mathematical operators for electronics
\newcommand{\F}[1]{\xmath{\mathcal{F}\left\{#1\right\}}} % fancy F for Fourier transform
\renewcommand{\L}[1]{\xmath{\mathcal{L}\left\{#1\right\}}} % fancy L for Laplace transform